% Options for packages loaded elsewhere
\PassOptionsToPackage{unicode}{hyperref}
\PassOptionsToPackage{hyphens}{url}
\documentclass[
]{article}
\usepackage{xcolor}
\usepackage{amsmath,amssymb}
\setcounter{secnumdepth}{-\maxdimen} % remove section numbering
\usepackage{iftex}
\ifPDFTeX
  \usepackage[T1]{fontenc}
  \usepackage[utf8]{inputenc}
  \usepackage{textcomp} % provide euro and other symbols
\else % if luatex or xetex
  \usepackage{unicode-math} % this also loads fontspec
  \defaultfontfeatures{Scale=MatchLowercase}
  \defaultfontfeatures[\rmfamily]{Ligatures=TeX,Scale=1}
\fi
\usepackage{lmodern}
\ifPDFTeX\else
  % xetex/luatex font selection
\fi
% Use upquote if available, for straight quotes in verbatim environments
\IfFileExists{upquote.sty}{\usepackage{upquote}}{}
\IfFileExists{microtype.sty}{% use microtype if available
  \usepackage[]{microtype}
  \UseMicrotypeSet[protrusion]{basicmath} % disable protrusion for tt fonts
}{}
\makeatletter
\@ifundefined{KOMAClassName}{% if non-KOMA class
  \IfFileExists{parskip.sty}{%
    \usepackage{parskip}
  }{% else
    \setlength{\parindent}{0pt}
    \setlength{\parskip}{6pt plus 2pt minus 1pt}}
}{% if KOMA class
  \KOMAoptions{parskip=half}}
\makeatother
\ifLuaTeX
  \usepackage{luacolor}
  \usepackage[soul]{lua-ul}
\else
  \usepackage{soul}
\fi
\setlength{\emergencystretch}{3em} % prevent overfull lines
\providecommand{\tightlist}{%
  \setlength{\itemsep}{0pt}\setlength{\parskip}{0pt}}
\usepackage{bookmark}
\IfFileExists{xurl.sty}{\usepackage{xurl}}{} % add URL line breaks if available
\urlstyle{same}
\hypersetup{
  hidelinks,
  pdfcreator={LaTeX via pandoc}}

\author{}
\date{}

\begin{document}

{
\setcounter{tocdepth}{3}
\tableofcontents
}
Paul Buy

The Pompeii Project

IRARAH answers

The open society and its enemies

A story from the Pompeii Project

\emph{\hl{How a professor disappeared in a Franciscan habit}}

\emph{\hl{The escape across the Tisza}}

Table of contents

\hyperref[vanished-without-a-trace]{\textbf{Vanished without a trace 2}}

\hyperref[the-monastery-in-simbach-am-inn]{\textbf{The monastery in Simbach am Inn 8}}

\hyperref[clothed-in-truth]{\textbf{Clothed in truth 9}}

\hyperref[escape-across-the-tisza]{\textbf{Escape across the Tisza 13}}

\hyperref[see-you-again-in-budapest]{\textbf{See you again in Budapest 18}}

\section{Vanished without a trace}\label{vanished-without-a-trace}

The rain fell in soft, even drops on the excavation site in the Pompeii Archaeological Park. The earth beneath the archaeologists\textquotesingle{} feet gradually transformed into a thick, muddy mass, while the rhythmic lapping of the water was the only sound that broke the silence. Thick, gray clouds hung over the ruins, so low that they seemed to swallow the surrounding hills. The world appeared shrouded in a damp veil; the ancient walls and unearthed relics seemed even more ephemeral, as if they might sink back into the ground at any moment.

Dr. Leonardo Moretti, the excavation director, stood leaning over a crumbling stone wall, his gray eyes fixed on the progress of the work. His weather-beaten face was serious, his thoughts drifting back to the centuries when these streets and buildings had still been teeming with the inhabitants of Pompeii. The past few days had yielded promising finds---fragments of inscriptions, well-preserved household objects. But today, a peculiar unease hung in the air, one that couldn\textquotesingle t be explained by the weather alone.

Suddenly, an assistant hurried towards him. His soaked clothes clung to his slender frame, mud splashing with every step. "Dr. Moretti, the inscription is almost uncovered. We need Martina Rossi for the assessment."

Moretti looked up from the wall. "Where is it? It should have been here long ago."

The assistant shrugged, a nervous twitch running across his face. "No one has seen her today. She wasn\textquotesingle t at breakfast either."

A strange feeling crept up inside Moretti---as if an invisible hand had closed around his stomach. Martina was reliable, a woman who took every appointment seriously. That she simply didn\textquotesingle t show up without a word was unusual. Too unusual to ignore. He looked at his watch. It was almost noon. The rain continued to patter against the stone floor.

``I'm going to check on her,'' he said, more to himself than to his assistant.

Moretti hurried to his car, which was parked at the edge of the excavation site. In his haste, he forgot his umbrella---a silly mistake. As he walked the short distance back, he noticed the drops from the tree leaves splashing onto his collar. The cold seeped through the fabric.

He revved the engine in the car. The windshield wipers glided across the windshield, their soft scraping mingling with the steady ticking of the clock on the dashboard. Moretti thought of Martina, of Julia. They were like sisters, inseparable -- even outside of work.

What had happened? His head was full of thoughts swirling around in a chaotic jumble. Had they hurt each other? Had there been an accident? He dismissed the thought. He would have heard about it.

He cast one last glance at the grounds, then turned into the narrow, rain-soaked street that led to her apartment.

He stopped in front of the small Italian house. The warm ochre facade with its crumbling plaster was familiar -- he had been there often. But on this rainy day, the house seemed different. The shutters rattled softly in the wind, rain dripped from the eaves. An invisible threat seemed to hang over the place.

Moretti knocked. No answer. He knocked louder---still nothing. He listened to the silence, hoping that any moment he would hear footsteps, that one of the women would open the door. But there was no sound.

He peered through a tilted window. His heart began to beat faster.

The apartment\textquotesingle s interior was in chaos. Clothes lay scattered on the beds as if hastily rummaged through. A half-packed suitcase stood askew in the hallway, its lid open. Papers and notes lay strewn across the kitchen table, as if someone had been hurriedly searching for something important. The scene had an unreal quality---but the disorder spoke volumes about a sudden, unprepared departure.

Moretti stepped back from the window. "That\textquotesingle s not like Martina," he thought. "She\textquotesingle s always so tidy."

He ran to his car, revved the engine, and sped off with spinning tires. Water splashed up from the roadside. He had to get help.

It was a gray, rainy morning when Moretti stepped through the heavy glass doors of the police station in Naples. The damp smell of the city clung to his clothes, and his wet shoes squeaked on the marble floor.

``I want to file a missing persons report,'' he said, his voice urgent and exhausted. ``Two of my colleagues---Martina Rossi and her mother, Julia Rossi---have been missing since yesterday evening. They were supposed to be at the excavations this morning. Their apartment\ldots'' He searched for words. ``It's chaos. As if they were in a hurry.''

The police officer at reception scrutinized him. "When did you last see the two of them?"

``Yesterday afternoon, at work. Everything was normal. They were near the new excavation site as usual. But after that, I didn\textquotesingle t hear from them again. They simply\ldots{} disappeared.'' Moretti\textquotesingle s voice broke slightly, a rare tremor he couldn\textquotesingle t suppress.

The police officer took down the information. "You say your apartment was in disarray? Were there any signs of violence?"

Moretti shook his head. ``No, nothing like that. But it wasn't normal -- the suitcases were half-packed, clothes were scattered everywhere, as if they were about to leave in a hurry.''

``We will take your statement and begin our investigation,'' the police officer said. His tone was reassuring, but Moretti sensed it wasn't enough.

He took a step back, looked around the station -- officials hurried past him, telephones rang. Moretti felt out of place. Here, in a world of order and regulations, he could only hope that the unease in his chest would soon give way to an answer. But deep down, he knew: This was only the beginning.

The police took the case seriously. Two detectives---an older man with graying hair, a younger man with a determined expression---took over the file. ``Two women, mother and daughter, missing since last night,'' the older one read. ``Apartment in chaos, no signs of a struggle.'' They exchanged a meaningful glance, then he swung open the car door.

They drove through the rain-soaked streets of Naples towards Pompeii. The sky remained heavy with clouds, and the drizzle collected on the windshield. In this region, where the shadows of the Camorra were ever-present, there were many reasons why two women could vanish without a trace.

As they reached the apartment building, the older man got out and pulled his coat tighter. The apartment door stood ajar---a detail that immediately caught his attention. ``This isn't good,'' he muttered. They entered cautiously. The air was stale and cold, the light dim, the silence unnatural.

Clothes were scattered across the beds. The half-packed suitcase lay in the hallway, its lid open, a lone shoe on the floor. Crumpled notes lay on the kitchen table. The younger man bent down and picked one up. "It looks like they were in the middle of preparing for a trip."

"That looks like a hasty departure," the older man muttered. "But no signs of a struggle." He opened the bathroom door---nothing out of the ordinary. Everything was normal, except for the mess.

The younger man raised the blinds. The dim daylight streamed into the room. "Perhaps they fled. There\textquotesingle s no indication of violence."

Back at the police station, the atmosphere was tense. The investigators were gathering their findings. Just another missing person case? There were too many unanswered questions.

"This is all taking place in an area where the Camorra has its fingers in the pie," the older man remarked. "It wouldn\textquotesingle t be the first time people have vanished without a trace." His colleague nodded.

The case was handed over to the public prosecutor\textquotesingle s office -- and to the Direzione Distrettuale Antimafia.

An oppressive silence hung in the prosecutor\textquotesingle s office. The prosecutor, a middle-aged man with deep wrinkles around his eyes, leafed through the initial reports. "A case like this near Pompeii could be connected to the Camorra," he murmured. "We have to follow every lead---bank accounts, phone activity, contacts. No detail is too insignificant."

The investigators stood tensely. One of the younger officers stepped forward. "The workplace has confirmed the missing person report. The colleagues at the excavation site are extremely concerned. I suggest we conduct another thorough investigation there."

The prosecutor nodded. "Good. And I want the DDA involved." He thought for a moment. "And keep an eye out for anything related to historical artifacts. There are many valuable excavations in the area---organized crime is interested in them."

Meanwhile, ARS, the artificial intelligence secretly operating for I.R.A.R.A.H., was preparing its next digital deception. Its algorithms worked swiftly and efficiently within the networks of the police and airlines. Flight and travel records were manipulated, bookings canceled, and passenger lists falsified. It now appeared as if Martina and Julia had never left the city. The AI \hspace{0pt}\hspace{0pt}covered its tracks so thoroughly that even experienced investigators were trapped in a thicket of false leads.

But the authorities didn\textquotesingle t give up. The suspicion that the two historians might have discovered something that was meant to be kept hidden was too concrete. During a renewed search of the apartment, they found a business card -- Michael Phillips, professor at the Gregorian University in Rome.

The investigators examined the card thoughtfully. "Who is this man? Why did they have his business card?" His proximity to the Vatican piqued their interest. A connection to Rome\textquotesingle s religious and academic elite---everything suddenly seemed important.

A team was assembled and sent to Rome.

Michael Phillips sat in his office at the Gregoriana University. The bookshelves stretched from floor to ceiling. The light from his desk lamp cast long shadows across the wooden table. For days, he had felt the tension mounting around him---like a net slowly but inexorably tightening. The investigation in Naples had intensified. Martina and Julia were moving into the spotlight---and with them, he.

An encrypted message from ARS had arrived: ``The air traffic control records have been deleted. No further traces.'' It was a relief---but fleeting. The pressure mounted with every minute. He knew that the slightest mistake could jeopardize everything.

As dusk fell over Rome, Michael was still sitting at his desk. The sounds of the city drifted in through the open window---the hum of the streetlights, the distant hum of engines. A somber crescendo that made his thoughts race faster and faster. He knew the investigators would be arriving soon.

Then the knocking on the door.

A loud knock broke the silence. Michael had been expecting this moment, had prepared himself -- but the churning in his stomach wouldn\textquotesingle t go away. With a deep breath, he stood up, forced himself to calm down, and opened the door.

Two men in dark suits stood before him. ``Dr. Phillips? We're from the Polizia di Stato. It's about the disappearance of Martina Rossi and Julia Rossi. May we come in?''

Michael nodded and stepped aside. They entered his office. He sensed the subdued tension in their movements---as if they were registering every wrinkle on his face, every involuntary gesture. He led them to the round table. The serious police officer sat opposite him; the other remained at the edge of the room.

``Do you know Martina Rossi and Julia Rossi very well?'' the police officer began.

``Yes. I have worked with them on various projects -- academically and within the framework of InSim.''

The police officer nodded curtly. "Martina\textquotesingle s employer has filed a missing person report. She was last seen in your vicinity. Can you tell us what happened on the day she disappeared?"

Michael paused for a moment. ``We met before their departure for Pompeii. They were going back to Italy for a workshop. Everything seemed normal.''

The men exchanged a quick glance. "Normal? There were no signs that anything was wrong?"

Michael shook his head. "Nothing I noticed." But inside, he was fighting. At that moment, the invisible interface of ARS buzzed softly in his ear: "They are reviewing aircraft recordings. We have deleted them. Stay calm."

"And what do you know about the InSim Mercedes accident near Pompeii? Two witnesses claim the vehicle was being followed."

Michael felt sweat on his forehead, but he forced himself to remain calm. ARS whispered: "We\textquotesingle ve processed the surveillance data. You won\textquotesingle t find any evidence."

"All I know is that they were on their way to a conference. The accident came unexpectedly. I was in Rome at the time."

The police officer watched him. Then he pulled out a business card -- Michael\textquotesingle s own. "This card was found among the women\textquotesingle s personal belongings. Can you explain that?"

``Yes. I gave them to both of them in case they wanted to contact me with academic questions.''

A deep silence fell. Michael felt the tension, the investigators\textquotesingle{} waiting. He kept his composure. ARS\textquotesingle s voice remained constant in his ear.

After what seemed like an eternity, the second investigator stood up and went to the window. "You don\textquotesingle t know anything about her current whereabouts?"

"Unfortunately not. I\textquotesingle m worried about her too."

He came back and leaned over the table. "If you\textquotesingle re hiding something from us, Dr. Phillips, we\textquotesingle ll find out."

Michael smiled thinly. "I understand. You can contact me anytime."

The men got up and left. The door clicked shut. Michael sank into his chair.

"They\textquotesingle re gone. The traces have been erased. Martina and Julia are safe," ARS whispered.

Michael closed his eyes and took a deep breath. But the worry remained -- buried deep inside him. They were all just a small step away from being exposed.

\section{The monastery in Simbach am Inn}\label{the-monastery-in-simbach-am-inn}

The monastery of the Congregatio Jesu in Simbach am Inn was almost deserted. Its impending closure made it the perfect hiding place.

Martina and Julia sat at a long wooden table in a former classroom. Michael\textquotesingle s doppelganger leaned against the wall. The faces of ARS, Michael Phillips, and an I.R.A.R.A.H. operator appeared on the laptop screen.

ARS\textquotesingle{} voice filled the room. "Welcome. The situation requires utmost discretion."

Michael Phillips, joining from Rome, appeared tense. "Professor Neumann in Kassel has received several threatening letters. His public lecture tomorrow could lead to an escalation. We must get him to safety before it\textquotesingle s too late."

Security footage of demonstrators outside the University of Kassel appeared on the screen -- angry faces, red signs.

``Professor Neumann is a respected scholar,'' explained the I.R.A.R.A.H. operator. ``He is known for his critical stance towards postmodern movements and transhumanism. Our task is to remove him from the university and place him in a safe location.''

Julia asked: "What if they discover us on the way?"

Michael replied: ``ARS has identified several alternative routes. If we don\textquotesingle t reach the chapel as planned, there are underground passages dating back to the Second World War. I.R.A.R.A.H. has prepared a secure hiding place in a chapel on the outskirts of Kassel.''

The doppelganger stepped away from the wall. "We should bring the professor here to the monastery after the rescue. The resolution is delayed. A safe place -- for a few days."

ARS nodded. "I will increase surveillance."

Michael Phillips raised his voice. ``Remember -- this rescue is more than an escape. It's about freedom of thought. About the right to seek the truth.''

The briefing ended. The mission began.

\section{Clothed in truth}\label{clothed-in-truth}

The morning was cool and foggy when Julia, Martina, and Michael\textquotesingle s doppelgangers arrived on the campus of the University of Kassel. A leaden haze hung over the buildings, the tense atmosphere almost palpable. The day was not beginning like any other, and they felt it to their very core. Even in the distance, they saw the first protest signs raised against the gloomy sky: "Against fascism and enemies of science!" and "Down with the reactionary!" echoed from a group of demonstrators gathered in front of the main building.

"This won\textquotesingle t be easy," Julia whispered, looking around with a worried expression. Her eyes scanned the crowd, crammed together in an angry mosaic of faces and signs.

``We need to hurry,'' said Michael's doppelganger, straightening his shoulders, which stood like a shield against the approaching cold. ``The professor is waiting for us. The faster we get him away from here, the less chance they have of recognizing us.''

Martina nodded and cast a searching glance at the grey clouds covering the sky, as if summoning an unwritten omen. "I.R.A.R.A.H. has prepared everything. Let\textquotesingle s not waste any time."

They slipped unnoticed into one of the annexes, their steps hurried but controlled. The sound of the protests outside grew louder, but inside the university building it was silent. The contrast felt surreal.

In the sparsely furnished office, they found Dr. Tobias Neumann, who was hurriedly packing his belongings into a worn leather bag. He was a middle-aged man with sharp features and an expression of exhaustion and determination in his eyes, which were deep in their sockets like two shadows in a dark alley.

When the team entered, he looked up and breathed a sigh of relief. "I\textquotesingle ve been waiting for you," he said, putting down his bag. "The situation outside is escalating. The protesters are particularly aggressive today." His voice was a rough thread that cut through the tense air.

``We have everything prepared,'' Michael's double said calmly, taking a brown Franciscan habit from his bag. ``Here is your religious habit. As soon as you put it on, you are officially Brother Timothy -- a Franciscan on his way to the USA.''

Martina handed him a forged identity card. "I.R.A.R.A.H. has ensured that you have a new identity. Your religious name and your new existence are secure." The words weighed heavily on the professor\textquotesingle s shoulders, who felt the weight of his situation.

Dr. Neumann stared at the habit before standing up resolutely. "This is crazy," he muttered as he wrapped himself in the brown robe. "But I have no other choice."

As soon as he finished, they led him through the empty corridors of the university building and out into the open. An inconspicuous delivery van was parked in front of the door, but the journey there was not without risk. The protests outside the university had grown louder; the crowd seemed agitated in anticipation of the professor\textquotesingle s upcoming lecture.

``We have to get through this,'' Julia said, glancing at the crowd. ``They mustn't realize it's you. We don't have much time.''

``I will go ahead with him,'' Michael's double said resolutely. ``They should believe we are a group of Franciscans on a pilgrimage.''

They moved slowly towards the crowd. The demonstrators barely noticed them -- until a shout rang out from the crowd: ``That's him! The reactionary professor!''

The protesters\textquotesingle{} eyes turned on her. For a moment, the situation seemed about to escalate. But Michael\textquotesingle s doppelganger quickly pushed Dr. Neumann into the van. With a dull thud, the door closed, and they drove off, angry shouts echoing behind them.

The highway to Frankfurt was quiet, but the tension was still palpable, like a rope being stretched to its limit. The professor leaned back and sighed heavily. "Discussions used to be possible," he said thoughtfully. "You could have a different opinion without being insulted. Today, all I see is anger and ignorance. I hope things will be different in the USA."

"Maybe," said Michael\textquotesingle s doppelganger. "Maybe not. But you\textquotesingle ll be safer there -- for now."

Arriving at Frankfurt Airport, the team escorted the professor through security. Every step had to be meticulously planned, as a single mistake could jeopardize their entire operation. Julia glanced over her shoulder as the professor handed his documents to the security officers. The officers were tired and efficient. They scanned the IDs, nodded, and waved them through.

Only one person hesitated -- a young officer with glasses, who looked at Neumann\textquotesingle s ID card longer than necessary.

"Brother Timothy? You are traveling to Chicago?"

``A pilgrimage,'' Neumann said calmly. ``To the Franciscans in the USA.''

The official nodded. "Have a good trip, brother."

They said goodbye at the gate. "Once you\textquotesingle re in the US, you\textquotesingle re safe," Julia said, placing a hand on the professor\textquotesingle s shoulder. "I.R.A.R.A.H. will take care of everything else."

Dr. Neumann nodded, a trace of gratitude in his eyes. "Without your help, I would have been lost," he said softly. "I owe you my life."

They watched him as he walked through the gate and the flight to Chicago was called. They took one last look at the man they had rescued before he disappeared behind the glass doors.

Upon arriving in the US, the professor was met by a member of the Franciscan community and taken to the Franciscan University of Steubenville. The team accompanied him on the long journey to the university, whose calm and peaceful atmosphere contrasted sharply with the chaotic conditions in Germany. The trees surrounding the campus seemed like silent sentinels, protecting the new arrivals.

``Welcome, Brother Timothy,'' one of the brothers greeted him with a warm smile. ``Your reputation precedes you. We are delighted to welcome you to our community.''

``It will be an honor,'' the professor replied, nodding slightly. The weight of his new identity felt both liberating and oppressive.

The day after his arrival, the professor delivered his inaugural lecture to a gathered group of Franciscans and students. The team sat at the back of the hall and listened attentively as the professor carefully chose his words. The hall was filled with a quiet anticipation; the air seemed to vibrate as he spoke.

``In a time when man is surrendering his autonomy to technology,'' he began, ``we must return to the values \hspace{0pt}\hspace{0pt}of the Enlightenment and rationality. It is not technology that will save us, but critical thinking. We must defend humanism against postmodern trends.''

The words echoed through the room. A collective nod went through the rows; the students seemed to sense the spark of hope that resided within the professor.

After the lecture, Martina, Julia, and Michael\textquotesingle s doppelgangers met with an I.R.A.R.A.H. agent in one of the university\textquotesingle s back rooms. The agent was a slim man with a serious face.

``It was risky,'' he said, his voice calm and controlled. ``But we have information that some of the demonstrators in Germany aim to infiltrate the Franciscan community as well. They might be targeting Dr. Neumann.''

``What can we do?'' Julia asked worriedly.

"We must ensure the professor\textquotesingle s safety and strengthen the connection to the community. He has become a target. The question is not whether they will try to find him, but when."

``That means we need to strengthen our defense,'' Martina added. ``He's not just a professor. He's a symbol.''

The agent nodded. "In the coming weeks, it will be crucial to secure our communications and keep a close eye on developments within the community. It will take time to smooth out the ripples that the events in Germany bring here."

In the following days, they prepared. Dr. Neumann gave further lectures and slowly found peace with his new identity. Once a week, the lecture hall filled with students who listened intently to his ideas. Here, discussion was possible again.

One evening, after a particularly inspiring talk, he sat at the table with his new brothers. "I never thought I\textquotesingle d feel so alive," he said, raising his glass. "To freedom of spirit!"

``To freedom of the spirit!'' the others shouted in unison, and a feeling of community filled the room.

For the first time in a long time, Dr. Neumann felt free. But he also knew that the team that had rescued him would face new dangers -- on the Ukrainian-Romanian border, on the Tisza River, in the darkness of night.

\section{Escape across the Tisza}\label{escape-across-the-tisza}

It was still early morning when the team gathered in the briefing room of the Franciscan University in Steubenville. The first rays of sunlight broke through the morning light, bathing the room in a soft gold, but the atmosphere was anything but relaxed. The air was filled with a mixture of concentration and nervousness -- everyone sensed that what lay before them could have far-reaching consequences.

A large map hung on the wall. The planned route through Romania and along the Tisza River was marked in bright red. The map lines seemed to pulsate, as if the route itself were breathing. On a large screen, the Zoom conference flickered -- the faces of ARS, Michael Phillips, and Agent Novak were visible.

Agent Novak, the IRAH operations commander, was already in front of the camera, his gaze steady and focused. ``Good morning, everyone,'' he began, his voice conveying both authority and concern. ``Our mission is clear: We must bring two men safely out of Ukraine---a Ukrainian pacifist and a Russian deserter. The Tisza River is the last barrier we must cross before returning to Romania. This region is under heavy surveillance, and we have only a small window of time.''

A murmur went through the room. Everyone knew: The lives of these men lay in their hands.

ARS\textquotesingle s calming, almost human-sounding voice filled the room. "The escape route was carefully planned. We recorded the positions of the border patrols and selected the safest point for crossing the Tisza River. Communication during the mission will take place via encrypted channels. Michael will remain in contact at all times via satellite phone."

Michael Phillips, who was acting as team leader, nodded in agreement. His eyes reflected a mixture of confidence and concern. "Remember---these men\textquotesingle s lives are in our hands. One wrong move could mean we\textquotesingle re exposed. So stay calm and focused. It\textquotesingle s crucial that we work together as a team."

He looked into the faces of his colleagues. He noticed the determination -- but also the uncertainty that hung like a shadow over the room.

"We have prepared everything, and I trust in your abilities. If we stick together, we can overcome this challenge."

A brief nod, an encouraging smile here and there -- the tension slowly began to ease.

``Now for the details,'' Michael continued, turning back to the map. ``We'll meet in a small village on the banks of the Tisza River. The contact man will be waiting for us there; he'll lead us to the men. The escape must be quick and quiet---no light signals, no noise. Each of us has a role. Any questions?''

A few hands went up. Michael answered the questions with clear, precise answers. As he spoke the last word, he sensed the team was ready.

"Okay, we don\textquotesingle t have much time. Gather your equipment and meet in the garage area in half an hour. Everyone knows what to do."

The team stood up. Michael stayed a moment longer, studying the map in front of him. The red lines seemed to pulsate.

``ARS,'' he murmured. ``Are there any risks we've overlooked?''

"All relevant information is included in the analysis. The probabilities are favorable -- as long as we adhere to the set timeframe and do not deviate from the plan."

Michael took a deep breath. The mission began.

The team boarded a flight from New York to Bucharest. In first class, where they had taken their seats, the atmosphere was one of tense anticipation. Captain Lukas Berger, an experienced ocean-going captain with a deep understanding of the dangers of escape, leaned slightly forward.

"The current on the Tisza is unpredictable," he explained with a serious look. "We have to be fast and quiet. The boat\textquotesingle s engine is muted, but any movement can attract attention. The area is under strict surveillance."

Julia listened attentively. "Each of us must do our part perfectly. Our mistakes must not endanger the freedom of others."

In another row, Michael\textquotesingle s doppelganger and Dr. Neumann---the professor they had already brought to safety---were talking. Their conversation was quiet, almost hushed. They were discussing the increasing surveillance through technology and the creeping erosion of personal freedom.

``Transhumanism and technocracy go hand in hand,'' Neumann murmured. ``They are trying to stifle discourse through control. Technology is used as a tool in this process.''

The doppelganger nodded vigorously. "That\textquotesingle s why we must show that freedom is something different from technological superiority. That is the mission of I.R.A.R.A.H. Our values \hspace{0pt}\hspace{0pt}are inextricably linked to human dignity."

After landing in Bucharest, they felt the pressure of the impending task. The journey to Sighetu Marmatiei, a border town on the Tisza River, was long -- almost nine hours of driving through picturesque landscapes, but one marked by the remnants of past conflicts.

ARS took control of communications. "You are approaching the river," the reassuring voice said over the encrypted radio. "Stick to the planned route. I am monitoring patrol movements in real time."

The landscape was melancholic -- vast fields, old, weathered villages that whispered stories of hope and despair.

Arriving at the agreed location, they parked the vehicle in a secluded, shady spot. The dense undergrowth protected them from prying eyes.

The faint sound of the river filled the air. They moved on foot through the thick undergrowth. With every step, the tension grew. Finally, they spotted the boat provided by I.R.A.R.A.H. -- hidden among tall trees, its engine muted, ready to launch.

The sky was covered with threatening, dark clouds -- a harbinger of their mission.

They climbed into the boat. Captain Berger took the helm. The engines started -- a rush of adrenaline broke the silence.

Berger navigated through the cold, dark waters of the Tisza River. The current was stronger than he had expected, but his experience kept him calm and focused.

``We are approaching the rendezvous point,'' ARS reported. ``The men are hidden in a wooded area on the Ukrainian shore. We only have a few minutes to find them.''

Everyone stared intently at the horizon. Then -- two figures emerged from the bushes. Their faces were etched with exhaustion, fear, and hope.

``Quick, come on board!'' Julia shouted, hastily helping them into the boat.

As soon as the men were safely on board, Berger pushed back towards the Romanian shore with full force.

Then a flash of light appeared on the horizon -- a patrol was visible in the distance.

"Hurry up!" Berger shouted. The engine roared, the boat shot through the water. The double held a cloth cloak over his side to hide it from the spotlights.

"Stay calm!" whispered Julia. "Everything will be alright. We\textquotesingle re on the right track."

They finally reached the Romanian shore. Hastily, they pulled the boat into the thicket. A local IRAH contact was already waiting.

``We'll sink it so there are no traces,'' he explained quietly, his eyes nervously scanning the surroundings. ``ARS showed us a safe escape route. But we have to hurry. The patrols could be here any minute.''

They helped the rescued men out of the boat and led them through the thicket to a small, secluded Jesuit monastery. The Jesuits had already made arrangements -- new identities, a safe place.

An older priest with a calm gaze and gentle smile greeted them. ``You risked a great deal to bring these men here. They are now safe. We will take care of them.''

Julia felt a wave of relief. The pressure on her shoulders lifted.

After bringing the rescued people to safety, the team set off for Hungary. They chose a remote route across the border to avoid being detected.

Michael contacted Julia via satellite phone. ``You did it. I.R.A.R.A.H. has confirmed that no trace of your presence was left behind. Good work.''

``Thank you,'' Julia replied. ``But we know this is just the beginning.''

They reached a small town in northeastern Hungary, where they were able to regroup.

Julia was preparing to begin her work as a social worker and psychological counselor. Martina wanted to pursue a career in archaeology. But deep down, they both knew: the time of peace would be short. Their escape across the Tisza River was only a small victory in the fight against oppression; the waves of change would not cease.

In the following days, news of renewed repression in Ukraine and Russia became increasingly urgent. Agent Novak contacted them with new instructions.

"The situation is escalating," he said. "We have reports of an imminent large-scale military exercise at the border. We must remain vigilant and be ready to act immediately."

``This is bothering us,'' Julia murmured. ``But we are ready. Always ready.''

The dangerous mission on the Tisza River was complete. But the danger was not over. Their journey had made them stronger; the determination in their hearts burned brighter than ever.

Together they would accept the next mission that I.R.A.R.A.H. entrusted to them -- no matter how challenging it might be.

\section{See you again in Budapest}\label{see-you-again-in-budapest}

The cool morning air enveloped Michael as he stood before the entrance of the Collegium Germanicum in Rome. The stone building, with its ancient walls, seemed to bid him a silent farewell. He felt the weight of the years he had spent there and the memories stored within its walls. The sun broke through the clouds, bathing the facades in a soft, golden light. With one last glance at his familiar surroundings, he let the heavy wooden door close behind him.

Outside, Maria waited for him, wrapped in a light cloak that fluttered gently in the morning breeze. Her gaze was calm, yet there was an expression in her eyes that made Michael pause for a moment. It was the look of a woman who had much to say---but who kept her words trapped in silence.

"So, the time has come," she said, her voice soft but tinged with a melancholy that made the air around her heavy. "You\textquotesingle re leaving and leaving everything behind."

Michael nodded slowly. "Yes, it\textquotesingle s time. There\textquotesingle s a lot to do in Budapest." He paused briefly and looked deeply into her eyes. "I hope you know that I still think of you---and everything we shared."

A faint smile flickered across Maria\textquotesingle s lips, but her eyes betrayed a deeper story. "Take care of yourself, Michael. It\textquotesingle s good to know you haven\textquotesingle t forgotten the family."

For a moment, it seemed as if there was more to her words than she let on. Michael knew that the hint about the past and the doppelganger\textquotesingle s identity hung between them like an unspoken shadow. He inclined his head, and without another word, he turned to go to the taxi that would take him to the airport.

-\/-\/-

The flight to Budapest was uneventful. As he gazed out the window at the passing landscape below -- the Alps, then the vast Hungarian Plain -- his thoughts revolved around what lay ahead and what he had left behind in Rome. The thought of Maria, of the unspoken words that hung between them, pushed its way to the forefront.

As the plane landed at Budapest airport, a familiar tingling of anticipation coursed through his body. It wasn\textquotesingle t just a new task that awaited him -- but also the opportunity to finally gain clarity.

A black car was waiting for him. The drive through the city took him past the Danube and the magnificent buildings that glowed in the golden evening light. The water sparkled in the twilight, seemingly reflecting the past and the future. He was on his way to the small apartment that Julia and Martina now called home.

As the car pulled up in front of the old building, Michael took a deep breath. Ancient trees framed the entrance, and the familiar sounds of the city filled the air. He pressed the doorbell and waited.

-\/-\/-

The door opened, and Martina greeted him with a warm smile that eased some of the tension from his shoulders. "Michael, it\textquotesingle s so good to see you. Come in, we\textquotesingle ve been waiting for you."

The living room had a cozy atmosphere. The aroma of freshly brewed coffee filled the room, and cakes and pastries, resembling miniature works of art, were laid out on the coffee table. Julia came in from the kitchen, carrying a tray, which she set down with a radiant smile.

``Finally you're here,'' she said, looking at Michael with an expression of relief. ``Sit down, have a coffee. We have a lot to talk about.''

Michael sat down on one of the old sofas, upholstered in a lovingly worn fabric. The room\textquotesingle s coziness gave him a feeling of home that he had long missed.

Shortly afterwards, the doppelganger also entered the room. He seemed relaxed -- but also a little nervous -- as he sat down opposite Michael.

``Welcome to Budapest,'' he said with a slight smile that betrayed both openness and uncertainty. ``It's nice to have the whole `family' together.''

The word "family" sounded strangely familiar to Michael. The expression on his doppelgänger\textquotesingle s face seemed, for a brief moment, to betray a deeper connection. Michael returned the smile---while a thought flashed through his mind: Was it really possible that this young man was his son?

``Thank you,'' said Michael, taking a sip of his coffee, which was warm and aromatic. ``It feels good to be here.''

-\/-\/-

They chatted about trivial things -- the city, work, life in Budapest. But between the conversations, an unspoken question hung in the air, a tension that neither Julia nor Martina seemed able to resolve. The doppelganger occasionally glanced at Michael, his eyes filled with intense interest, as if seeking confirmation of something Michael hadn\textquotesingle t yet voiced.

Dinner was relaxed. The conversation was light and full of laughter. But as they finally sat down on the balcony in the twilight, a kind of unspoken understanding hung in the air. Michael knew it was time to seek answers---and that he might already have them right in front of him.

``Sometimes,'' he began softly, as the first stars twinkled in the sky, ``life leads us down unexpected paths that we only understand later.'' He looked at his doppelganger and noticed that he was listening intently to his words. ``And sometimes we meet people who show us that there are more connections than we initially believe.''

The doppelganger said nothing. But he nodded.

Night fell over Budapest, and the city lights twinkled like tiny sparks, casting memories into the darkness. The glances they exchanged spoke volumes as the evening stillness enveloped them. Michael knew it would take time to speak the truth completely---but for now, it was enough that they were together.

The family had taken on a new dimension -- one he hadn\textquotesingle t expected, but perhaps had always hoped for.

The gentle murmur of the Danube could be heard in the distance. Michael knew: This was only the beginning of a long and exciting journey.

\end{document}
